\section*{Part V: Structured Assessments: Making Important Capabilities Observable}

\subsection*{43. Design the Assessment Portfolio Before Choosing an Interview Format}

A selection battery should be built backward from the role's critical competencies. Imagine a 4-column planning sheet: \textit{capability}, \textit{method}, \textit{evidence produced}, and \textit{decision importance}. The founder might choose a work sample for applied technical judgment, a situational interview for integrity under pressure, a reference for repeated reliability, and a short eligibility screen for scheduling constraints. There is no benefit in conducting an elaborate personality test if the selection criteria concern technical review accuracy and turnaround.

An effective portfolio covers the most important work while keeping candidate burden tolerable. A common mistake is to distribute competencies by interviewer (``you assess culture, you assess technical, you assess motivation'') without defining what any of those words means. Another is to use 5 interviews to measure the same social performance and then call the consensus scientific.

Create 1 assessment blueprint for each distinct role. It should state what each stage measures, why that characteristic matters, how it will be scored, what threshold applies, and how results feed the decision. The blueprint becomes the institution's protection against improvisation.

\subsection*{44. Structure Is What Makes an Interview an Assessment}

A structured interview uses predetermined job-related questions, standardized prompts or permitted probes, defined rating anchors, and a consistent process for recording and evaluating responses. Structure reduces the opportunity for irrelevant impressions to dominate. Professional guidance from the U.S. Office of Personnel Management notes that higher interview structure is generally associated with better validity and inter-rater agreement [R6].

Structure does not mean treating candidates as machines. Interviewers can listen attentively, offer clarification, and invite evidence. The key is that \textit{the information elicited and the scoring standards remain comparable}. If 1 candidate receives a difficult technical scenario and another is asked about hobbies, the resulting scores do not measure the same thing. If 1 candidate receives extensive hints and another receives none, comparisons become unreliable.

Plan the interview in advance. State how long it lasts, which topics it covers, who scores each competency, and what follow-up questions are permissible. Give interviewers a rubric with examples of weak, adequate, and excellent responses. Require evaluators to record observed evidence before the overall rating.

\subsection*{45. Behavioral Questions Investigate What Candidates Actually Did}

A behavioral interview asks for a specific past episode: \textit{Tell me about a time you identified a serious problem that others had missed. What information did you have, what did you do, and what happened?} It seeks actual conduct rather than self-description. Strong answers identify an event, the candidate's own decisions, constraints, tradeoffs, result, and subsequent learning.

Interviewers should distinguish ``we'' from ``I'' without assuming collaboration is undesirable. Ask what the team accomplished and which part the candidate personally performed. Probe for a decision that could have gone differently. Ask what evidence was available at the time. A candidate who says the project succeeded but cannot explain the failure modes may have been peripheral: or may simply need clearer prompts. Avoid converting uncertainty into an accusation.

Behavioral interviews are informative when the candidate has relevant experience, but past opportunity is unevenly distributed. A candidate who has not managed a major production incident cannot provide such a story. For early-career roles, provide situational questions or a simulation alongside historical examples so that selection does not simply reward prior access to prestigious assignments.

\subsection*{46. Situational Questions Investigate Judgment Under a Shared Scenario}

A situational interview gives each candidate an analogous problem and asks how they would respond. For a reviewer: \textit{A learner challenges a low score, and your preliminary review suggests that the rubric itself may be ambiguous. The final report is due tomorrow. What do you do?} This can reveal how candidates balance standards, fairness, urgency, and uncertainty.

A strong rubric would reward checking the original evidence, acknowledging the possibility of an institutional error, documenting an interim decision, obtaining a 2nd review where necessary, communicating with the learner, and correcting systemic causes. A weak response might ignore the complaint, invent certainty, or change the grade simply to eliminate conflict. The scoring standard is tied to job behavior, not the elegance of the candidate's speech.

Situational responses reveal a plan, not proof that the candidate will execute it. A candidate may know the socially appropriate answer without having the discipline to act. That is why scenarios and historical evidence complement 1 another. For high-stakes conduct, past references and supervised work may be especially informative.

\subsection*{47. Build Behaviorally Anchored Scoring Scales}

A rating anchor translates a vague judgment into observable criteria. Consider the competency \textbf{diagnosing faulty technical reasoning} in a project-review role. A 5-point scale can be defined as shown in Table~\ref{tab:bars_scale}.

\begin{table}[htbp]
\centering
\small
\caption{Behaviorally Anchored Rating Scale (BARS): Diagnosing Faulty Technical Reasoning}
\label{tab:bars_scale}
\begin{tabularx}{\columnwidth}{l X}
\toprule
\textbf{Score} & \textbf{Behavioral Anchor} \\
\midrule
1: Unsafe & Misses a material error or endorses a fundamentally invalid conclusion \\
2: Weak & Notices a symptom but cannot explain its causal or statistical mechanism \\
3: Meets threshold & Identifies the main issue and gives a technically correct, actionable fix \\
4: Strong & Also distinguishes primary from secondary defects and anticipates likely learner misunderstandings \\
5: Exceptional & Adds a reproducible check, explains uncertainty, and gives a concise route to independent verification \\
\bottomrule
\end{tabularx}
\end{table}

The scale should be tested using examples. Do not assume that 5 labels alone create 5 meaningful levels. If evaluators cannot distinguish score 3 from score 4, collapse categories or improve anchors. Scoring also needs instructions for partial answers and uncertain evidence.

Behavioral anchors protect against the halo effect: being articulate, credentialed, or agreeable should not automatically inflate a candidate's technical score. They also provide a basis for calibration and respectful feedback.

\subsection*{48. Independent Scoring Should Precede Group Discussion}

In a panel, 1 respected interviewer can influence everyone else's judgment. A confident opening comment (``this candidate is clearly brilliant'') may cause others to reinterpret ambiguous evidence favorably. A negative comment can do the reverse. To limit this, evaluators should record their observations and assign preliminary ratings \textbf{before} discussing the candidate. Significant discrepancies should then be investigated by referring to the evidence and rubric.

For instance, 2 assessors may give the same demonstration scores of 2 and 5 because 1 evaluated lecture fluency while the other evaluated whether the candidate corrected a learner misconception. That is not a useful compromise to average at 3.5. It reveals a disagreement about the construct or scoring standard. Resolve that disagreement before the next candidate.

A panel can still make a holistic final judgment after individual evidence is collected. The important safeguard is that individual assessments are not erased by social consensus, especially when some interviewers have greater organizational status.
